% 三角錐数に 1 を加えた平方数の決定
% 組版: LuaLaTeX（lualatex を 2〜3 回）。
\RequirePackage{fix-cm}
\documentclass[a4paper,11pt]{ltjsarticle}

\usepackage{amsmath,amssymb,amsthm,mathtools}
\usepackage{booktabs,array}
\usepackage{enumitem}
\usepackage[hidelinks]{hyperref}
\usepackage{newunicodechar}
\newunicodechar{ℕ}{\ensuremath{\mathbb{N}}}

\setlist{nosep,leftmargin=2em}
\setlist[enumerate]{label=(\arabic*)}
\newcolumntype{L}[1]{>{\raggedright\arraybackslash}p{#1}}

\newtheoremstyle{jplain}{\topsep}{\topsep}{\normalfont}{}{\bfseries}{}{1em}{}
\theoremstyle{jplain}
\newtheorem{theorem}{定理}[section]
\newtheorem{proposition}[theorem]{命題}
\newtheorem{lemma}[theorem]{補題}
\newtheorem{corollary}[theorem]{系}
\newtheorem{definition}[theorem]{定義}
\newtheorem{remark}[theorem]{注意}
\renewcommand{\proofname}{証明}
\makeatletter
\renewenvironment{proof}[1][\proofname]{\par\pushQED{\qed}\normalfont\topsep6\p@\@plus6\p@\relax
  \trivlist\item[\hskip\labelsep\bfseries #1]\ignorespaces}{\popQED\endtrivlist\@endpefalse}
\makeatother
\numberwithin{equation}{section}
\numberwithin{table}{section}

\newcommand{\Z}{\mathbb{Z}}
\newcommand{\Q}{\mathbb{Q}}
\newcommand{\N}{\mathbb{N}}
\newcommand{\F}{\mathbb{F}}
\newcommand{\Rr}{\mathbb{R}}
\newcommand{\Nm}{\operatorname{N}}
\newcommand{\NF}{\operatorname{N}_F}
\newcommand{\tsec}[2]{\texorpdfstring{#1}{#2}}
\newcommand{\lean}[1]{\texttt{#1}}

\title{三角錐数に $1$ を加えた平方数}
\author{roiban1344}
\date{2026年11月}

\begin{document}
\maketitle

\begin{abstract}
非負整数 $n$ について、三角錐数 $n(n+1)(n+2)/6$ に $1$ を加えた数が平方数になるのは、$n=0,5,8,22,24,63,13775$ のときに限る。
この方程式は楕円曲線を定めるので、整点の決定には通常、楕円対数の一次形式の下界に基づく計算が使われる。
本稿は、楕円曲線の理論を使わない証明を与える。
証明の中心は、二つの方程式を、相対ノルムが $11$ である一つの元にまとめることにある。
その元を $11$ で割って単数の冪に帰着し、指数を $2$ 進の一意性補題で決める。
必要な有限計算は、具体的な整数の恒等式と合同式に限られる。証明の全体は Lean~4 と Mathlib で形式化されている。
\end{abstract}

\tableofcontents

%======================================================================
\section{序}
%======================================================================

\subsection{主定理}

以下、$\N=\{0,1,2,\dots\}$ とし、$T(n)=n(n+1)(n+2)/6$ を三角錐数とする（OEIS A000292）。

\begin{theorem}\label{thm:main}
$n\in\N$ について、ある $m\in\N$ で $T(n)+1=m^2$ となるのは、
\[
n\in\{0,\ 5,\ 8,\ 22,\ 24,\ 63,\ 13775\}
\]
のときであり、そのときに限る。対応する $m$ は順に $1,6,11,45,51,209,660099$ である。
\end{theorem}

$m$ の列は OEIS A216269「$k^2-1$ が三角錐数となる $k$」として、この七項で登録されている。

\subsection{楕円曲線との関係}

方程式 $T(n)+1=m^2$ を $6$ 倍すると $6m^2=n(n+1)(n+2)+6$ となる。
$X=6n$, $Y=36m$ と置けば、これは
\[
Y^2=X^3+18X^2+72X+1296
\]
となり、楕円曲線の整点の問題である。
したがって、解が有限個であることは Siegel の定理から従う。
しかし Siegel の定理は解の大きさの上界を与えない。
そのため、整点をすべて求めるには通常、楕円対数の方法を用いる~\cite{ST94,GPZ94}。
この方法では、Mordell--Weil 群の基底を求め、楕円対数の一次形式の下界から係数の上界を得て、最後に格子簡約で残りを調べる。

本稿の証明は、このいずれも使わない。
使うのは次のものだけである。
\begin{itemize}
\item 階数 $4$ の自由 $\Z$ 加群に具体的な乗法を入れた五つの環での計算。
\item 実数の不等式。
\item 整数の合同式と整除。
\end{itemize}
類数、単数定理、イデアルの一意分解、$p$ 進解析は使わない。
単数については、具体的な一つの元がすべての相対単数を生成することを、実数の不等式で直接示す（第~\ref{sec:units} 節）。
また指数の決定には、$p$ 進の零点の理論の代わりに、$2$ の冪による整除を繰り返す議論を用いる（第~\ref{sec:dyadic} 節）。

\subsection{証明の構成}

証明は次の順に進む。
\begin{enumerate}
\item 方程式を五つの連立方程式 $n+3=dr^2$, $n^2+2=ey^2$ に分ける（第~\ref{sec:branches} 節）。
\item 第二式の解を、二つの整数による平方の形で表す（第~\ref{sec:parameters} 節）。
\item 二つの式をまとめて、相対ノルムが $11$ の元 $\xi$ を作る（第~\ref{sec:norm11} 節）。
\item $11$ による割り算で、$\xi$ を有限個の初期元と相対単数の積に書く（第~\ref{sec:division} 節）。
\item 相対単数が一つの元 $\rho$ の冪と符号で尽きることを示す（第~\ref{sec:units} 節）。
\item 負の符号と、不要な初期元を除く（第~\ref{sec:signs} 節）。
\item $2$ 進の一意性補題で、$\rho$ の指数を決める（第~\ref{sec:dyadic} 節と第~\ref{sec:exponents} 節）。
\item 決まった元の座標から $n$ を読む（第~\ref{sec:conclusion} 節）。
\end{enumerate}
最後の段では、元 $\xi$ の第 $0$ 座標がはじめから $n+3$ に等しいので、他の量を経由せずに $n$ を読める。

\subsection{形式化}

すべての段は Lean~4 と Mathlib で形式化されている。
最終定理は付録~\ref{app:lean} に示す形をしており、Lean のカーネルが検査する。
依存する公理は Lean の標準の三つ（\lean{propext}, \lean{Classical.choice}, \lean{Quot.sound}）だけである。
\lean{sorry} や \lean{native\_decide} は使わない。
有限の計算は、カーネルが \lean{decide} で確かめる。
本文の各段と Lean のファイルの対応は、付録~\ref{app:lean} の表にまとめた。

\subsection{記号}

$e$ を平方数でない正の整数とし、$\omega=\sqrt e$ と書く。
$\Z[\omega]=\{a+b\omega : a,b\in\Z\}$ のノルムを $\NF(a+b\omega)=a^2-eb^2$ とする。
$\Z[\omega]$ の元 $\alpha$ の共役 $a-b\omega$ を $\alpha'$ と書く。

%======================================================================
\section{五つの連立方程式}\label{sec:branches}
%======================================================================

$T(n)+1=m^2$ を $6$ 倍して整理すると
\begin{equation}\label{eq:shift}
6m^2=(n+3)(n^2+2)
\end{equation}
を得る。

\begin{proposition}\label{prop:branches}
$(n,m)$ が \eqref{eq:shift} の自然数解ならば、正の整数 $r,y$ と、表~\ref{tab:branches} の五つの組 $(d,e)$ のいずれかについて
\begin{equation}\label{eq:pair}
n+3=dr^2,\qquad n^2+2=ey^2
\end{equation}
が成り立つ。
\end{proposition}

\begin{table}[htbp]
\centering
\caption{五つの場合}\label{tab:branches}
\begin{tabular}{c|ccccc}
\toprule
$d$ & $1$ & $2$ & $3$ & $11$ & $66$\\
$e$ & $6$ & $3$ & $2$ & $66$ & $11$\\
\bottomrule
\end{tabular}
\end{table}

\begin{proof}
$n+3=dr^2$ と書き、$d$ を平方因子を持たない正の整数とする。
まず、$n+3$ と $n^2+2$ の公約数は $11$ を割る。実際、$(n^2+2)-(n-3)(n+3)=11$ である。
次に、\eqref{eq:shift} の両辺で素因数の指数の偶奇を比べると、$d$ は $6(n^2+2)$ を割る。
したがって、$d$ を割る素数 $p$ が $6$ を割らなければ、$p$ は $n+3$ と $n^2+2$ の公約数であり、$p=11$ となる。
よって $d\mid66$ である。

また、\eqref{eq:shift} から $6d(n^2+2)$ は平方数である。
そこで $6d=k^2e$（$e$ は平方因子を持たない）と書くと、$e(n^2+2)$ も平方数であり、$n^2+2=ey^2$ を得る。
$d=1,2,3,6,11,22,33,66$ に対して $e=6,3,2,1,66,33,22,11$ である。

残るのは $d=6,22,33$ の除外である。
$d=6,22$ では $e=1,33$ となる。どちらも $e\equiv1\pmod4$ なので $n^2+2\equiv y^2\pmod4$ となるが、平方数は法 $4$ で $0$ か $1$ なので、これは不可能である。
$d=33$ では $n+3=33r^2$ から $3\mid n$ である。すると $n^2+2=22y^2$ から $22y^2\equiv2\pmod3$、すなわち $y^2\equiv2\pmod3$ となり、これも不可能である。
\end{proof}

逆に、\eqref{eq:pair} の解は \eqref{eq:shift} の解を与える。実際、$de=6k'^2$（$k'=1,1,1,11,11$）なので、$m=k'ry$ とすればよい。

%======================================================================
\section{二次の平方表示}\label{sec:parameters}
%======================================================================

この節では、第二式 $n^2+2=ey^2$ の解を、二つの整数 $a,b$ の二次式で表す。
Pell 方程式の解の一般論は使わず、具体的な降下でパラメータを構成する。

\subsection{降下}

$e\ne2$ の四つの場合、$e=a_0^2+2$（$a_0=2,1,8,3$）と書ける。
$c_0=a_0^2+1$ と置くと、$c_0+a_0\omega$ はノルム $1$ を持つ。

\begin{lemma}\label{lem:descent}
$n\ge0$, $y>1$ が $n^2+2=ey^2$ を満たすとき、
\[
N=c_0n-a_0ey,\qquad Z=c_0y-a_0n
\]
は $0\le N<n$, $Z>0$, $N^2+2=eZ^2$ を満たす。さらに $n=c_0N+a_0eZ$, $y=c_0Z+a_0N$ である。
\end{lemma}

\begin{proof}
$N+Z\omega=(n+y\omega)(c_0-a_0\omega)$ なので、ノルムは保たれる。最後の二式は、$c_0+a_0\omega$ を掛けて戻したものである。
大小関係は次の順に示す。$e=c_0+1$ に注意する。
\begin{enumerate}
\item $a_0y<n$。そうでなければ $n^2\le a_0^2y^2=ey^2-2y^2<ey^2-2=n^2$ となる（$y>1$ を使う）。
\item $Z>0$。恒等式 $(c_0y)^2-(a_0n)^2=y^2+2a_0^2>0$ による。
\item $|N|<n$。$N>-n$ は $(c_0+1)n>a_0ey$、すなわち $n>a_0y$ と同値である。$N<n$ は $a_0n<ey$ と同値であり、(2) から $a_0n<c_0y<ey$ である。
\item $a_0Z\le|N|$。(1) と同じ議論を $(|N|,Z)$ に使う（$Z\ge1$ による）。
\item $N\ge0$。$N<0$ と仮定すると、$n=c_0N+a_0eZ=a_0eZ-c_0|N|\le a_0eZ-c_0a_0Z=a_0Z\le|N|<n$ となり、矛盾する。
\end{enumerate}
\end{proof}

この降下は $y=1$ で止まり、そのとき $n=a_0$ である。

\begin{proposition}\label{prop:parameters}
$e=A^2+2$ とし、$Q_A(a,b)=a^2-2b^2+2Aab$ と置く。
$n\ge0$, $y>0$ が $n^2+2=ey^2$ を満たすならば、$A=a_0$ または $A=-a_0$ と整数 $a,b$ があって
\begin{equation}\label{eq:param}
Q_A(a,b)=1,\qquad n=A(a^2-2b^2)-4ab
\end{equation}
が成り立つ。
\end{proposition}

\begin{proof}
$n$ についての帰納法で示す。$y=1$ のときは $n=a_0$ であり、$(A,a,b)=(a_0,1,0)$ でよい。
$y>1$ のときは、補題~\ref{lem:descent} の $(N,Z)$ は $0\le N<n$ を満たすので、帰納法の仮定から表示を持つ。そこから次の更新で $(n,y)$ の表示を得る。
\[
(a_0,a,b)\mapsto(-a_0,\ a,\ b-a_0a),\qquad(-a_0,a,b)\mapsto(a_0,\ a-2a_0b,\ b).
\]
\eqref{eq:param} が保たれることは、展開して確かめられる。
\end{proof}

\subsection{符号の選択}

\eqref{eq:param} の二式から $a^2-2b^2=1-2Aab$ を代入すると、
\begin{equation}\label{eq:offset}
n=A-2(A^2+2)ab
\end{equation}
を得る。この式と第一式 $n+3=dr^2$ を合わせると、各場合で $A$ の符号の一方が除かれる。
\begin{itemize}
\item $d=1$, $A=2$：$Q_2(a,b)=1$ を法 $4$ で見ると $b$ は偶数である。すると $r^2=n+3=5-12ab\equiv5\pmod8$ となり、不可能である。
\item $d=2$, $A=1$：$2r^2=4-6ab$ から $r^2\equiv2\pmod3$ となり、不可能である。
\item $d=11$, $A=-8$：$11r^2=-5-132ab$ から $11\mid5$ となり、不可能である。
\item $d=66$, $A=3$：$66r^2=6-22ab$ から $11\mid6$ となり、不可能である。
\end{itemize}
したがって $A=-2,\,-1,\,8,\,-3$（$d=1,2,11,66$）である。

\subsection{\tsec{$e=2$}{e=2} の場合}

$d=3$, $e=2$ では $n+3=3r^2$ から $3\mid n$ であり、さらに $n$ は偶数である（$n^2+2=2y^2$ より）。
$n=2N$ と書くと $y^2-2N^2=1$ である。
法 $4$ で見ると $y$ は奇数、$N$ は偶数である。
互いに素な連続整数 $(y-1)/2$, $(y+1)/2$ の積が $N^2/2$ に等しいので、各因子は平方数または平方数の $2$ 倍である。
ここからノルム $\pm1$ の元 $a+b\sqrt2$ で $y+N\sqrt2=(a+b\sqrt2)^2$ となるものを得る。
このとき $n=2N=4ab$ である。
もしノルムが $-1$ なら、$a^2+b^2\equiv2\pmod3$ から $3\nmid ab$ となり、$3\mid n$ に反する。
したがってノルムは $1$ であり、$b$ の符号を変えれば \eqref{eq:param} の $A=0$ の表示になる。

\subsection{平方表示}

\begin{proposition}\label{prop:square}
$n^2+2=ey^2$, $y>0$ であり、\eqref{eq:param} が成り立つならば、$y=a^2+2b^2$ であり、
\[
n+y\omega=(A+\omega)\bigl(a+(A-\omega)b\bigr)^2
\]
が成り立つ。
\end{proposition}

\begin{proof}
恒等式
\[
\bigl(A(a^2-2b^2)-4ab\bigr)^2+2Q_A(a,b)^2=(A^2+2)(a^2+2b^2)^2
\]
に \eqref{eq:param} を代入すると、$n^2+2=e(a^2+2b^2)^2$ を得る。$y>0$ と比べて $y=a^2+2b^2$ である。後半は展開による。
\end{proof}

後の環に合わせるため、既知の単数を一度掛けて固定元を取り替える。
結果は $w=n+y\omega=\beta\eta^2$, $\eta=\eta_0+\eta_1\omega$ の形になり、各場合の $\beta,\eta_0,\eta_1$ は表~\ref{tab:beta} のとおりである。
表の各行は $a,b$ についての二次の恒等式であり、展開で確かめられる。

\begin{table}[htbp]
\centering
\caption{平方表示 $n+y\omega=\beta(\eta_0+\eta_1\omega)^2$}\label{tab:beta}
\begin{tabular}{ccccc}
\toprule
$d$ & $A$ & $\beta$ & $\eta_0$ & $\eta_1$\\
\midrule
$1$ & $-2$ & $22+9\omega$ & $5a+2b$ & $-2a-b$\\
$2$ & $-1$ & $5+3\omega$ & $2a+b$ & $-a-b$\\
$3$ & $0$ & $\omega$ & $a$ & $-b$\\
$11$ & $8$ & $8+\omega$ & $a+8b$ & $-b$\\
$66$ & $-3$ & $63+19\omega$ & $10a+3b$ & $-3a-b$\\
\bottomrule
\end{tabular}
\end{table}

%======================================================================
\section{相対ノルムが \tsec{$11$}{11} の元}\label{sec:norm11}
%======================================================================

\subsection{五つの環}

各場合に、$\Z^4$ に乗法を入れた可換環 $R=R_d$ を一つ定める。
$R$ は $\Z[\omega]$ を部分環として含み、$\theta^2=\Delta\in\Z[\omega]$ を満たす元 $\theta$ を持つ。
$d\ne11$ では $R=\Z[\omega]\oplus\Z[\omega]\theta$ であり、$d=11$ では $R=\Z[\omega]\oplus(\Z+\Z\,\omega/11)\theta$ である。
$\Delta$ と座標の意味は表~\ref{tab:rings} のとおりである。

\begin{table}[htbp]
\centering
\caption{五つの環。$x=(x_0,x_1,x_2,x_3)$ の表す元と $\theta^2=\Delta$}\label{tab:rings}
\begin{tabular}{cccc}
\toprule
$d$ & $e$ & $x$ の表す元 & $\Delta=\theta^2$\\
\midrule
$1$ & $6$ & $(x_0+x_1\omega)+(x_2+x_3\omega)\theta$ & $44+18\omega$\\
$2$ & $3$ & 同上 & $20+12\omega$\\
$3$ & $2$ & 同上 & $6\omega$\\
$11$ & $66$ & $(x_0+x_1\omega)+(x_2+x_3\omega/11)\theta$ & $176+22\omega$\\
$66$ & $11$ & $(x_0+x_1\omega)+(x_2+x_3\omega)\theta$ & $756+228\omega$\\
\bottomrule
\end{tabular}
\end{table}

$d=1,2,3,11$ では $\Delta=2d\beta$ である。$d=66$ では $2d\beta=11\Delta=(\omega\theta)^2$ である。
以下、$s=\theta$（$d\ne66$）、$s=\omega\theta$（$d=66$）と置く。どの場合も $s^2=2d\beta$ である。

\begin{definition}
\begin{enumerate}
\item $R$ の\textbf{相対共役} $\tau$ を、$\theta\mapsto-\theta$ で定める環の自己同型とする。座標では $\tau(x)=(x_0,x_1,-x_2,-x_3)$ である。
\item $x\in R$ の\textbf{相対ノルム}を $\Nm(x)=x\,\tau(x)$ とする。これは $\Z[\omega]$ に属する。
\item $\Nm(x)=1$ となる $x$ を\textbf{相対単数}と呼ぶ。相対単数 $x$ は可逆で、$x^{-1}=\tau(x)$ である。
\item $x$ の最初の二座標を $P(x)=x_0$, $Q(x)=x_1$ と書き、$\pi(x)=x_0+x_1\omega\in\Z[\omega]$ と置く。
\end{enumerate}
\end{definition}

各環の乗法表は、表~\ref{tab:rings} の関係から決まる。
Lean では、これらを $\Z^4$ 上の具体的な乗法として定義し、可換環の公理を確かめている。
数体の整数環であることは仮定しない。

\subsection{相対ノルムが \tsec{$11$}{11} の元}

\begin{proposition}\label{prop:xi}
\eqref{eq:pair} の解に対し、$\xi=w+3+r\eta s\in R$ と置く。ここで $w=n+y\omega$ である。このとき
\begin{equation}\label{eq:norm11}
\Nm(\xi)=11
\end{equation}
であり、さらに
\begin{equation}\label{eq:f}
f(\xi)=0,\qquad P(\xi)=n+3=dr^2
\end{equation}
が成り立つ。ここで $f(x)=(P(x)-3)^2-eQ(x)^2+2$ である。
\end{proposition}

\begin{proof}
$\NF(w)=n^2-ey^2=-2$ に注意すると、
\begin{align*}
\xi\,\tau(\xi)&=(w+3)^2-r^2\eta^2s^2=(w+3)^2-2dr^2\beta\eta^2\\
&=(w+3)^2-2(n+3)w=w^2-2nw+9
\end{align*}
となる。さらに $w+w'=2n$, $ww'=-2$ なので $w^2-2nw=-ww'=2$ であり、$\Nm(\xi)=11$ を得る。
$\pi(\xi)=w+3$ なので、$P(\xi)=n+3$, $Q(\xi)=y$ であり、$f(\xi)=n^2-ey^2+2=0$ である。
\end{proof}

ここで $11=3^2+2$ である。二つの式 \eqref{eq:pair} が、一つの小さな相対ノルムにまとまる。
$\xi$ の座標は、$d=1,2,3$ で $(n+3,y,r\eta_0,r\eta_1)$、$d=11$ で $(n+3,y,r\eta_0,11r\eta_1)$、$d=66$ で $(n+3,y,11r\eta_1,r\eta_0)$ である。

関数 $f$ は最初の二座標だけで決まる。したがって $f(\tau(x))=f(x)$, $P(\tau(x))=P(x)$ である。

%======================================================================
\section{\tsec{$11$}{11} による割り算}\label{sec:division}
%======================================================================

表~\ref{tab:initial} の元を\textbf{初期元}と呼ぶ。どれも相対ノルムは $11$ である。

\begin{table}[htbp]
\centering
\caption{初期元と相対単数 $\rho$}\label{tab:initial}
\begin{tabular}{cccc}
\toprule
$d$ & $\xi_0$ & $\xi_1$ & $\rho$\\
\midrule
$1$ & $(1,1,5,-2)$ & — & $(9,4,2,0)$\\
$2$ & $(2,1,2,-1)$ & $(2,3,1,0)$ & $(3,2,1,0)$\\
$3$ & $(3,1,1,0)$ & — & $(17,12,6,4)$\\
$11$ & $(11,1,1,0)$ & — & $(33,4,2,2)$\\
$66$ & $(0,1,0,0)$ & — & $(19,6,1,0)$\\
\bottomrule
\end{tabular}
\end{table}

\begin{proposition}\label{prop:division}
$\Nm(\xi)=11$ ならば、ある初期元 $p$ と相対単数 $u$ があって、$\xi=u\,p$ または $\xi=u\,\tau(p)$ となる。
\end{proposition}

\begin{proof}
いずれの場合も、$\xi\,\tau(p)$ または $\xi\,p$ の四座標がすべて $11$ で割れることを示せばよい。
実際、そのとき $u=\xi\tau(p)/11$ と置けば $\Nm(u)=\Nm(\xi)\Nm(p)/121=1$ であり、$\xi\,\tau(p)\,p=11\xi$ から $\xi=up$ となる。

$d=1,3$ の場合、$e=6,2$ は法 $11$ で平方剰余でない。
したがって $\Z[\omega]/(11)$ は $121$ 元の体 $\F_{121}$ である。
$\xi_0=\alpha_0+\beta_0\theta$ と書くと、$\Nm(\xi_0)=\alpha_0^2-\Delta\beta_0^2=11$ であり、$\beta_0$ は $11$ と互いに素である。
したがって $\Delta$ は $\F_{121}$ で平方根 $\pm\delta$（$\delta=\alpha_0/\beta_0$）を持つ。
$\theta\mapsto\pm\delta$ で二つの環準同型 $\chi_\pm:R\to\F_{121}$ が定まり、$\chi_-=\chi_+\circ\tau$ である。
$\chi_+(\xi)\chi_-(\xi)=\chi_+(\Nm(\xi))=0$ なので、$\xi$ はどちらかで $0$ になる。
$\xi_0$ と $\tau(\xi_0)$ はそれぞれ別の側で $0$ になるので、$\xi$ と同じ側で $0$ になる方を $p$ とすれば、$\xi\tau(p)$ は両方で $0$ になる。
両方の核の共通部分は $11R$ である。

$d=2$ の場合、$e=3$ は法 $11$ で平方剰余であり、$\Z[\omega]/(11)\cong\F_{11}\times\F_{11}$ となる。
二つの成分のそれぞれで、$\theta$ の値の選び方が二通りある（$\Delta$ の平方根の存在は上と同様に $\xi_0$ から分かる）。
そこで、$\xi$ が各成分でどちらの値で $0$ になるかの組を、$\xi$ の\textbf{型}と呼ぶ。型は四通りある。
$\xi_0,\xi_1,\tau(\xi_0),\tau(\xi_1)$ がこの四つの型を一つずつ実現するので、上と同じ議論で、$\xi$ と同じ型を持つ $p$ について $\xi\tau(p)\in11R$ となる。

$d=11$ の場合、相対ノルムを法 $11$ で読むと $x_0^2-96x_3^2\equiv0$ となる。
$96\equiv8$ は法 $11$ で平方剰余でないので、$11\mid x_0$, $11\mid x_3$ である。
ここから $\xi\tau(\xi_0)\in11R$ が従う。

$d=66$ の場合も同様に、$x_0^2-756x_2^2\equiv0\pmod{11}$ と $756\equiv8$ から $11\mid x_0$, $11\mid x_2$ である。
ここから $\xi\xi_0\in11R$ が従う。$\tau(\xi_0)=\xi_0$ に注意する。
\end{proof}

この割り算は、この特定の $11$ だけについて示している。類数やイデアルの一意分解は使わない。

%======================================================================
\section{相対単数}\label{sec:units}
%======================================================================

\begin{proposition}\label{prop:units}
各場合で、相対単数はすべて $\pm\rho^m$（$m\in\Z$）の形をしている。ここで $\rho$ は表~\ref{tab:initial} の元である。
\end{proposition}

$\rho$ が相対単数であることは直接の計算で確かめられる。示すべきは、他に相対単数がないことである。

\subsection{二つの実埋め込み}

$\Delta=\Delta_0+\Delta_1\omega$ について、$\sigma(\Delta)=\Delta_0+\Delta_1\sqrt e>0$, $\sigma'(\Delta)=\Delta_0-\Delta_1\sqrt e<0$ である。
$D_+=\sigma(\Delta)$, $D_-=-\sigma'(\Delta)$ と置く。
$\sigma:R\to\Rr$ を $\omega\mapsto\sqrt e$, $\theta\mapsto\sqrt{D_+}$ で定める。
また $\omega\mapsto-\sqrt e$ で定まる $\Z[\omega]$ の共役を $\alpha\mapsto\alpha'$ と書く。

相対単数を $v=\alpha+\gamma\theta$ と書く。ここで $\alpha\in\Z[\omega]$ であり、$\gamma\in\Z[\omega]$（$d=11$ では $\gamma\in\Z+\Z\,\omega/11$）である。
$\Nm(v)=\alpha^2-\Delta\gamma^2=1$ に共役を取ると、$\alpha'^2+D_-\gamma'^2=1$ を得る。したがって $|\gamma'|\le D_-^{-1/2}$ である。

\begin{lemma}[空隙]\label{lem:gap}
$\gamma\ne0$ ならば
\begin{equation}\label{eq:gap}
\bigl(|\sigma(v)|-|\sigma(v)|^{-1}\bigr)^2\ge4M,\qquad M=\nu^2D_+D_-
\end{equation}
が成り立つ。ここで $\nu=1$（$d\ne11$）、$\nu=3/11$（$d=11$）である。
\end{lemma}

\begin{proof}
$\sigma(\tau(v))=\sigma(v)^{-1}$ なので、$\sigma(v)-\sigma(v)^{-1}=2\sigma(\gamma)\sqrt{D_+}$ である。
$\gamma\ne0$ なら $|\sigma(\gamma)\gamma'|=|\NF(\gamma)|\ge\nu$ である。
実際、$d\ne11$ では $\NF(\gamma)$ は零でない整数である。
$d=11$ では $\gamma=c+k\omega/11$ と書くと $\NF(\gamma)=(11c^2-6k^2)/11$ であり、分子の絶対値が $1,2$ になることは法 $3$ と法 $11$ で除かれる。
したがって
\begin{equation}\label{eq:gamma}
\sigma(\gamma)^2\ge\nu^2/\gamma'^2\ge\nu^2D_-
\end{equation}
であり、\eqref{eq:gap} を得る。
\end{proof}

$4M$ と $\Lambda=\sigma(\rho)$ の値は表~\ref{tab:gap} のとおりである。$\Lambda$ については、表の上界と $\Lambda>1$ だけを使う。

\begin{table}[htbp]
\centering
\caption{空隙の定数と $\Lambda=\sigma(\rho)$ の上界}\label{tab:gap}
\begin{tabular}{cccccc}
\toprule
$d$ & $1$ & $2$ & $3$ & $11$ & $66$\\
\midrule
$4M$ & $32$ & $128$ & $288$ & $288$ & $1152$\\
$\Lambda<$ & $40$ & $14$ & $100$ & $142$ & $84$\\
\bottomrule
\end{tabular}
\end{table}

\subsection{縮約}

\begin{proof}[命題~\ref{prop:units} の証明]
$\Lambda>1$ なので、任意の相対単数 $v$ に $\rho$ の整数冪を掛けて、$1\le a=|\sigma(v)|<\Lambda$ とできる。
$\gamma=0$ なら $\alpha^2=1$ であり、$\alpha\in\Z[\omega]$ から $v=\pm1$ となる。
$\gamma\ne0$ と仮定して矛盾を導く。

$v'=\rho\,\tau(v)$ と置くと、$b=|\sigma(v')|=\Lambda/a$ であり、$b>1$, $ab=\Lambda$ である。
$v'$ の $\theta$ 成分も零でない（零なら $v=\pm\rho$ となり $a=\Lambda$ に反する）。
補題~\ref{lem:gap} を $v$ と $v'$ に使い、$a,b\ge1$ で成り立つ $(a-a^{-1})^2\le a^2$ と合わせると、$a^2\ge4M$, $b^2\ge4M$ を得る。
したがって $\Lambda=ab\ge4M$ である。
$d=2,3,11,66$ では、表~\ref{tab:gap} から $\Lambda<4M$ なので矛盾する。

$d=1$ では $\Lambda>32$ なので、もう一段の議論が要る。
$b^2\ge32$ から $a^2=\Lambda^2/b^2<50$ である。
$v$ の $\Z[\omega]$ 成分 $\alpha$ について、$\sigma(\alpha)=(\sigma(v)+\sigma(v)^{-1})/2$ なので $\sigma(\alpha)^2<14$ である。
また $\alpha'^2\le1$ である。
$(\sigma(\alpha)+2\alpha')^2\ge0$ から
\[
(\sigma(\alpha)-\alpha')^2\le\tfrac32\sigma(\alpha)^2+3\alpha'^2<24
\]
を得る。
ところが $\alpha=h+c\sqrt6$ と書けば、左辺は $24c^2$ である。したがって $c=0$ である。
このとき $\sigma(\alpha)^2=\alpha'^2\le1$ となる。一方、$\sigma(\alpha)^2=1+D_+\sigma(\gamma)^2$ と \eqref{eq:gamma} から $\sigma(\alpha)^2\ge1+D_+D_-=9$ が従い、矛盾する。
\end{proof}

この議論は一般の単数定理を使わない。必要なのは、実数での縮約と、表~\ref{tab:gap} の有理数による上界だけである。

%======================================================================
\section{符号と第二の初期元}\label{sec:signs}
%======================================================================

前二節から、$\xi=\pm p\rho^m$（$p$ は初期元またはその相対共役、$m\in\Z$）と書ける。
$\tau$ は $f$ と第 $0$ 座標を保ち、$\tau(\rho)=\rho^{-1}$ を満たす。
したがって $\tau(\xi_i)\rho^m=\tau(\xi_i\rho^{-m})$ の形の元は、$\xi_i\rho^{-m}$ に置き換えてよい。
以下、$\xi=\pm\xi_i\rho^m$ を調べる。

\subsection{非負性}

\begin{lemma}\label{lem:positive}
各初期元 $\xi_i$ とすべての $m\in\Z$ について、$P(\xi_i\rho^m)\ge0$ である。
\end{lemma}

\begin{proof}
五つの環の乗法表の係数はすべて非負である。したがって、座標がすべて非負の元の全体 $C$ は積について閉じている。$\rho$ も $C$ に属する。

各初期元について、$\xi_i\rho^2$ と $\tau(\xi_i)\rho^2$ が $C$ に属することを直接確かめる。
すると $m\ge2$ では $\xi_i\rho^m=(\xi_i\rho^2)\rho^{m-2}\in C$ である。
$m\le-2$ では $\tau(\xi_i\rho^m)=\tau(\xi_i)\rho^{-m}\in C$ であり、$\tau$ は第 $0$ 座標を保つ。
残る $m=-1,0,1$ は直接確かめる。
\end{proof}

$P(\xi)=n+3>0$ なので、$\xi=-\xi_i\rho^m$ の形は起こらない。

\subsection{\tsec{$d=2$}{d=2} の第二の初期元}

\begin{lemma}\label{lem:extra}
$d=2$ で $\xi=\xi_1\rho^m$ となることはない。
\end{lemma}

\begin{proof}
$\rho^2=(41,24,6,4)$ を掛けても、最初の二座標は法 $8$ で変わらない。このことは任意の元について乗法表から確かめられる。
$m$ が偶数なら、$\xi_1\rho^m$ の最初の二座標は法 $8$ で $(2,3)$ であり、$f\equiv8\pmod{16}$ となって $f(\xi)=0$ に反する。
$m$ が奇数なら、第 $0$ 座標は法 $8$ で $4$ である。ところが $P(\xi)=2r^2$ であり、平方数は法 $4$ で $2$ にならないので、矛盾する。
\end{proof}

以上で、調べるべき元は $\xi_0\rho^m$（$m\in\Z$）だけになった。

%======================================================================
\section{\tsec{$2$}{2} 進の一意性補題}\label{sec:dyadic}
%======================================================================

この節では、$f(y\,W^i)=0$ から $i=0$ を導く一般的な補題を示す。
ここで $y$ は $f$ の零点であり、$W$ は $1$ に $2$ 進的に近い可逆元である。
$p$ 進解析の零点定理は使わず、$2$ の冪による整除を繰り返して示す。

以下、この補題を適用する $f$ の零点 $y$ を\textbf{基点}と呼ぶ。
また $W=1+2qU$ と書いたときの $U$ を、$W$ の\textbf{方向}と呼ぶ。

\subsection{展開}

$y,z\in R$, $c\in\Z$ について
\begin{equation}\label{eq:expand}
f(y+cz)=f(y)+c\lambda_y(z)+c^2H(z)
\end{equation}
が成り立つ。ここで
\[
\lambda_y(z)=2\bigl\{(P(y)-3)P(z)-eQ(y)Q(z)\bigr\},\qquad H(z)=P(z)^2-eQ(z)^2
\]
である。$\lambda_y$ は $z$ について $\Z$ 線形であり、$P(z),Q(z)$ だけで決まる。

基点 $y$ の座標を $(y_0,y_1,y_2,y_3)$ と書く。$f(y)=0$ を展開すると
\begin{equation}\label{eq:root}
y_0^2-ey_1^2=6y_0-11
\end{equation}
を得る。

\subsection{冪の展開}

\begin{lemma}\label{lem:power}
$t\ge0$ とし、$W=1+2^{t+1}U$ が可逆であるとする。このとき任意の $k\ge0$ と $\ell\in\Z$ について、次が成り立つ。
\begin{enumerate}
\item ある $E\in R$ があって $W^{2^k\ell}=1+2^{k+t+1}\bigl\{\ell U+2^tU^2E\bigr\}$ となる。
\item $t\ge1$ ならば、ある $b\in\Z$ と $E\in R$ があって $W^{2^k\ell}=1+2^{k+t+1}\bigl\{\ell U+2^tU^2(b+2E)\bigr\}$ となる。
\end{enumerate}
\end{lemma}

\begin{proof}
まず $W^{2^k}$ を、$k$ についての帰納法で二乗を繰り返して展開する。
次に、$\ell\ge0$ では $W^{2^k}$ の冪を二項展開し、$\ell<0$ では逆元 $W^{-2^k}$ の同じ形の展開を使う。
(2) では、誤差の項を $U^2$ の倍数として保つだけでなく、その係数のうち整数の部分 $b$ を分けて保つ。
こうすると、二乗を繰り返しても $b+2E$ の形が崩れない。
\end{proof}

以下、$q=2^t$ と置く。(2) の形は $W^{2^k\ell}=1+2^{k+1}q\{\ell U+qU^2(b+2E)\}$ と書ける。

\subsection{一意性補題}

\begin{lemma}[$2$ 進の一意性補題]\label{lem:unique}
$q=2^t$（$t\ge1$）とし、$W=1+2qU$ を可逆元、$y\in R$ とする。
$h\in\{2,4\}$ と $G\in R$ について $U^2=hG$ とし、$D=2qh^2$ と置く。次を仮定する。
\begin{enumerate}
\item $f(y)=0$。
\item $h\mid P(yU)$, $h\mid Q(yU)$。
\item $h=2$ であるか、または $h=4$ であって、$\chi(x)=P(x)+Q(x)\bmod2$ が積を保つ写像 $R\to\F_2$ であり、$\chi(G)=0$ を満たす。
\item $2h^2\mid\lambda_y(yU^2)$。
\item $\lambda_y(yU)\equiv qh^2\pmod D$。
\end{enumerate}
このとき、$i\in\Z$ について $f(yW^i)=0$ ならば $i=0$ である。
\end{lemma}

\begin{proof}
最初に、(1) と (3) から、任意の $E\in R$ について
\begin{equation}\label{eq:anyerror}
h^2\mid\lambda_y(yU^2E)
\end{equation}
が成り立つことを示す。
$h=2$ のときは、$\lambda_y$ がそもそも $2$ の倍数であり、$U^2=2G$ なので明らかである。
$h=4$ のときは、\eqref{eq:root} を法 $2$ で読む。平方は法 $2$ で元の数に等しいので、$P(y)-3\equiv eQ(y)\pmod2$ である。したがって
\[
\lambda_y(z)/2=(P(y)-3)P(z)-eQ(y)Q(z)\equiv(P(y)-3)\chi(z)\pmod2
\]
となる。$\chi(yGE)=\chi(y)\chi(G)\chi(E)=0$ から $4\mid\lambda_y(yGE)$ であり、$U^2=4G$ と合わせて \eqref{eq:anyerror} を得る。

次に、$k\ge0$ について「$2^k\mid i$ ならば $2^{k+1}\mid i$」を示す。
これが示されれば、$i$ はすべての $2$ の冪で割れるので、$i=0$ である。

$i=2^k\ell$ と書く。補題~\ref{lem:power} (2) から、$yW^i=y+cZ$, $c=2^{k+1}q$, $Z=y\{\ell U+qU^2(b+2E)\}$ となる。
\eqref{eq:expand} と $f(y)=0$ から、$0=f(yW^i)=c\{\lambda_y(Z)+cH(Z)\}$ であり、$\lambda_y(Z)+cH(Z)=0$ を得る。

ここで $\lambda_y(Z)=\ell\lambda_y(yU)+qb\,\lambda_y(yU^2)+2q\lambda_y(yU^2E)$ である。
\eqref{eq:anyerror} と (4)(5) から
\[
\lambda_y(Z)\equiv\ell qh^2\pmod{2qh^2}
\]
となる。
一方、(2) と $U^2=hG$ から $P(Z)$, $Q(Z)$ は $h$ の倍数であり、$H(Z)$ は $h^2$ の倍数である。$c$ は $2q$ の倍数なので、$cH(Z)\equiv0\pmod{2qh^2}$ である。
以上から $\ell qh^2\equiv0\pmod{2qh^2}$、すなわち $\ell$ は偶数であり、$2^{k+1}\mid i$ となる。
\end{proof}

以下では、仮定 (2)〜(4) が方向 $U$ の相対ノルムの算術から一般に従うことを示す。
個々の基点で確かめるのは、(1) と、$h=4$ のときの (5) だけである。

\subsection{方向の相対ノルム}

\begin{lemma}\label{lem:direction}
$W=1+4jU$（$j\in\Z$, $j\ne0$）が相対単数ならば、
\[
\pi(U)=-2j\Nm(U),\qquad U^2=-W\Nm(U)
\]
が成り立つ。
\end{lemma}

\begin{proof}
$\Nm(W)=1+4j(U+\tau U)+16j^2\Nm(U)=1$ であり、$U+\tau U=2\pi(U)$ である。これで第一式を得る。
第二式は $U^2=U(U+\tau U)-U\tau U=-4jU\Nm(U)-\Nm(U)$ による。
\end{proof}

$W=1+2qU$, $q\ge2$ はこの補題の形（$j=q/2$）をしている。

各環で、$\theta$ の係数どうしの積が $\pi$ に与える寄与は $h$ の倍数になる。ここで $h$ は次の値である。
\[
h=2\quad(d=1,3,11),\qquad h=4\quad(d=2,66).
\]
実際、$d\ne11$ では $\Delta$ の二係数が $h$ で割れる。$d=11$ では基底に $\omega\theta/11$ があるので、整数係数の式で直接確かめる。

\begin{corollary}\label{cor:width}
補題~\ref{lem:direction} の $W$ について、$\Nm(U)\in h\Z[\omega]$, $\pi(U)\in2h\Z[\omega]$, $U^2\in hR$ が成り立つ。
さらに任意の $y\in R$ について、$P(yU)$, $Q(yU)$ は $h$ の倍数である。
\end{corollary}

\begin{proof}
$\pi(U)=-2j\Nm(U)$ は偶数係数なので、$\pi(U)^2\in4\Z[\omega]$ である。
$\Nm(U)$ は $\pi(U)^2$ と $\theta$ の係数どうしの寄与の和なので、$\Nm(U)\in h\Z[\omega]$ である。すると $\pi(U)=-2j\Nm(U)\in2h\Z[\omega]$ となる。
$U^2=-W\Nm(U)$ から $U^2\in hR$ である。
最後に、$\pi(yU)$ は $\pi(y)\pi(U)$ と $\theta$ の係数どうしの寄与の和なので、$h$ の倍数である。
\end{proof}

これで補題~\ref{lem:unique} の (2) が従い、$G=U^2/h$ が定まる。

\subsection{仮定 (3)}

$h=4$ の二つの場合（$d=2,66$）では、$e=3,11$ は奇数であり、$\Delta$ は偶数係数である。
このとき $\chi(x)=x_0+x_1\bmod2$ は、$\omega\mapsto1$, $\theta\mapsto0$ という代入にほかならない。
この代入は $\omega^2=e\equiv1$, $\theta^2=\Delta\equiv0\pmod2$ と両立するので、$\chi$ は環準同型 $R\to\F_2$ である。

次に $\chi(G)=0$ を示す。$\chi$ は $\pi$ 部分だけで決まるので、$G=U^2/4$ の $\pi$ 部分を読めばよい。
$\pi(U)^2/4$ は、系~\ref{cor:width} から偶数係数である。
また $\theta$ の係数どうしの寄与は、$\Delta/4=5+3\omega$ または $189+57\omega$ の倍数である。
どちらも $\chi$ で $0$ になるので、$\chi(G)=0$ である。

\subsection{仮定 (4)}

$\Nm(U)=u+v\omega$ と置く。
補題~\ref{lem:direction} から $yU^2=-y\Nm(U)-4j(yU)\Nm(U)$ である。
後半の項は、系~\ref{cor:width} から $\pi$ 部分が $h^2$ の倍数なので、$\lambda_y$ の値は $8jh^2$ の倍数である。
前半の項については、\eqref{eq:root} を使うと
\[
\lambda_y(y\Nm(U))=2\bigl\{(3y_0-11)u-3ey_1v\bigr\}
\]
となる。したがって、括弧の中が $h^2$ の倍数であることを示せばよい。

$h=2$ の場合、$e$ は偶数であり、$u,v$ も偶数である。
また \eqref{eq:root} を法 $2$ で読むと $y_0^2\equiv1$ なので、$y_0$ は奇数である。
したがって $3y_0-11$ は偶数であり、括弧の中は $4$ の倍数である。

$h=4$ の場合は、次の二つの合同式を使う。$\varsigma=-1$（$d=2$）、$\varsigma=1$（$d=66$）と置く。
\begin{itemize}
\item 軌道の合同式：すべての $m\in\Z$ について、$y=\xi_0\rho^m$ は $y_0+\varsigma y_1\equiv1\pmod4$ を満たす。
実際、任意の $x$ について $(P+\varsigma Q)(x\rho)-(P+\varsigma Q)(x)$ は $4$ の倍数であり、$\xi_0$ での値は $1$ である。
\item 方向の合同式：$u-\varsigma v\equiv0\pmod{16}$ かつ $4\mid u$。
これは $\pi(U)\in8\Z[\omega]$ と、$\theta$ の係数を $c+k\omega$ と置いたときの $\Nm(U)$ の展開から従う。
\end{itemize}
$e\equiv-1\pmod4$ なので、$(3y_0-11)-3e\varsigma y_1=3(y_0+\varsigma y_1-1)-8-3\varsigma(e+1)y_1\equiv0\pmod4$ である。
これを $v\equiv\varsigma u\pmod{16}$, $4\mid u$ と合わせると、括弧の中は $16$ の倍数である。

\subsection{仮定 (5)：\tsec{$h=2$}{h=2} の場合}

$h=2$ の三つの場合では、(5) も一般に従う。
基点 $y$ の $\theta$ の係数を $\gamma$（$d\ne11$ では $y_2+y_3\omega$、$d=11$ では $y_2+y_3\omega/11$）と書く。
相対ノルムが $11$ であることは $\pi(y)^2-11=\Delta\gamma^2$ と書ける。両辺の $\NF$ を取ると、表~\ref{tab:J} の整数 $J_d$ について
\[
\NF(\pi(y)^2-11)=-8J_d^2
\]
となる。他方、展開により $f(y)f(-y)=\NF(\pi(y)^2-11)+8y_0^2$ である。
したがって $f(y)=0$ ならば $y_0^2=J_d^2$ である。

\begin{table}[htbp]
\centering
\caption{$\NF(\pi(y)^2-11)=-8J_d^2$ の $J_d$}\label{tab:J}
\begin{tabular}{ccc}
\toprule
$d=1$ & $d=3$ & $d=11$\\
\midrule
$y_2^2-6y_3^2$ & $3(y_2^2-2y_3^2)$ & $11y_2^2-6y_3^2$\\
\bottomrule
\end{tabular}
\end{table}

$h=2$ の三つの場合では $e\equiv2\pmod4$ である。
\eqref{eq:root} を法 $4$ で読むと $y_0,y_1$ は奇数であり、$J_d\equiv y_2^2\pmod2$ なので $y_2$ も奇数である。
また系~\ref{cor:width} から $P(U),Q(U)$ は $4$ の倍数である。
乗法表を法 $4$ で読むと、$P(yU)\equiv0$, $Q(yU)\equiv2y_2U_2\pmod4$ となる。ここで $U_2$ は $U$ の第 $2$ 座標である。
したがって $U_2$ が奇数ならば
\[
\lambda_y(yU)=2(y_0-3)P(yU)-2ey_1Q(yU)\equiv8\pmod{16}
\]
であり、これは $q=2$, $h=2$ の (5) にほかならない。
実際に使う三つの方向では $U_2=9,99,57$ で、いずれも奇数である。

$h=4$ の場合の (5) は、使う四つの基点で直接確かめる（表~\ref{tab:centres}）。

%======================================================================
\section{指数の決定}\label{sec:exponents}
%======================================================================

補題~\ref{lem:unique} を使う基点 $y$ と単数 $W$ は表~\ref{tab:centres} のとおりである。
基点はすべて $\xi_0\rho^{m_0}$ の形をしており、$f(y)=0$ は直接確かめる。

\begin{table}[htbp]
\centering
\caption{一意性補題を使う基点}\label{tab:centres}
\begin{tabular}{ccccccc}
\toprule
$d$ & 基点の指数 $m_0$ & $W$ & $q$ & $h$ & $D$ & (5) の確認\\
\midrule
$1$ & $0,\ -1$ & $\rho^2$ & $2$ & $2$ & $16$ & 一般\\
$2$ & $0,\ -4$ & $\rho^8$ & $4$ & $4$ & $128$ & 直接\\
$2$ & $-1$ & $\rho^4$ & $2$ & $4$ & $64$ & 直接\\
$3$ & $0,\ -1$ & $\rho^2$ & $2$ & $2$ & $16$ & 一般\\
$11$ & $0$ & $\rho^2$ & $2$ & $2$ & $16$ & 一般\\
$66$ & $1$ & $\rho^4$ & $2$ & $4$ & $64$ & 直接\\
\bottomrule
\end{tabular}
\end{table}

$W=1+2qU$ の形になることは、法 $4$ での $\rho$ の振る舞いから分かる。
$d=1,3,11$ では $\rho\equiv1\pmod2$ なので $\rho^2\equiv1\pmod4$ である。
$d=2,66$ では $\rho^2\equiv1+2\theta\pmod4$ であり、その二乗 $\rho^4$ が $1\pmod4$ に入る。
さらに $\rho^4=1+4U'$ なら $\rho^8=1+8(U'+2U'^2)$ なので、$d=2$ の $\rho^8$ は $q=4$ の形をしている。

以下、$f(\xi_0\rho^m)=0$ となる $m$ を決める。

\paragraph{$d=1,3$。}
$m$ が偶数なら $m=2a$、奇数なら $m=-1+2a$ と書き、それぞれ基点 $\xi_0$, $\xi_0\rho^{-1}$ で補題~\ref{lem:unique} を使う。
したがって $m=0$ または $m=-1$ である。

\paragraph{$d=11$。}
$\tau(\xi_0)=\xi_0\rho^{-1}$ なので、$\tau(\xi_0\rho^m)=\xi_0\rho^{-1-m}$ である。
$\tau$ は $f$ と第 $0$ 座標を保つので、$m$ を $-1-m$ に置き換えて $m$ は偶数としてよい。
基点 $\xi_0$ で補題~\ref{lem:unique} を使うと $m=0$ である。

\paragraph{$d=66$。}
まず、元の平方条件から $m$ が奇数であることを示す。この段では $f$ を使わない。
$\xi_0=\omega$ なので、$m=2i$ ならば $P(\xi)=11\,Q(\rho^{2i})=66r^2$、すなわち $Q(\rho^{2i})=6r^2$ である。
$\rho^2=1+2S$, $S=(756,228,19,6)$ と書き、$i=2^ka$ とする。
補題~\ref{lem:power} (1) を $W=\rho^2$, $t=0$ に使うと、$\rho^{2i}=1+cV$, $c=2^{k+1}$, $V=aS+S^2E$ を得る。
$\Nm(\rho^{2i})=1$ の $\omega$ 係数を比べ、$V$ の相対ノルムの $\omega$ 係数が法 $48$ で $12a^2$ に合同であることを使うと、
\[
r^2=-c^2(a^2+4g)
\]
となる整数 $g$ を得る。
$c^2\mid r^2$ から $r=cs$ と書け、$s^2+a^2\equiv0\pmod4$ となるので $a$ は偶数である。
この議論を $k=0,1,2,\dots$ と繰り返すと、$i$ はすべての $2$ の冪で割れるので $i=0$ となり、$Q(\rho^0)=0=6r^2$ が $r>0$ に反する。
したがって $m$ は奇数である。

次に、$\tau(\xi_0)=\xi_0$ なので $\tau(\xi_0\rho^m)=\xi_0\rho^{-m}$ である。
$m$ を $-m$ に置き換えて $m\equiv1\pmod4$ としてよい。
基点 $\xi_0\rho$ で補題~\ref{lem:unique} を使うと $m=1$ である。

\paragraph{$d=2$。}
この場合は、近い二つの根 $0,-4$ を持つので、一段の準備が要る。
軌道の元 $x=\xi_0\rho^m$ はすべて $2\mid x_0$, $x_1+2x_3\equiv3\pmod4$ を満たす。
この条件を満たす任意の $x$ について、乗法を展開すると
\[
f(x\rho^3)-3f(x\rho^2)+3f(x\rho)-f(x)\equiv0\pmod{128}
\]
が成り立つ。
最初の三つの値と、この式による正逆両方向の帰納法から、すべての $m\in\Z$ について
\begin{equation}\label{eq:entrance}
f(\xi_0\rho^m)\equiv32m(m+1)\pmod{128}
\end{equation}
を得る。
したがって $f(\xi_0\rho^m)=0$ ならば $m\equiv0$ または $m\equiv-1\pmod4$ である。
$m=4a$ のときは、$a$ の偶奇で分けて基点 $\xi_0$, $\xi_0\rho^{-4}$ と $W=\rho^8$ で補題~\ref{lem:unique} を使い、$m=0$ または $m=-4$ を得る。
$m=-1+4a$ のときは、基点 $\xi_0\rho^{-1}$ と $W=\rho^4$ で $m=-1$ を得る。

%======================================================================
\section{主定理の証明}\label{sec:conclusion}
%======================================================================

\begin{proof}[定理~\ref{thm:main} の証明]
以上で、$\xi$ は表~\ref{tab:values} の九つの元のいずれかである。
$P(\xi)=n+3$ なので、第 $0$ 座標から $n$ が決まる。

\begin{table}[htbp]
\centering
\caption{残る元と $n$}\label{tab:values}
\begin{tabular}{cccc}
\toprule
$d$ & $m$ & $P(\xi_0\rho^m)$ & $n$\\
\midrule
$1$ & $0$ & $1$ & $-2$\\
$1$ & $-1$ & $25$ & $22$\\
$2$ & $0$ & $2$ & $-1$\\
$2$ & $-4$ & $13778$ & $13775$\\
$2$ & $-1$ & $8$ & $5$\\
$3$ & $0$ & $3$ & $0$\\
$3$ & $-1$ & $27$ & $24$\\
$11$ & $0$ & $11$ & $8$\\
$66$ & $1$ & $66$ & $63$\\
\bottomrule
\end{tabular}
\end{table}

負の値 $n=-2,-1$ は $n\in\N$ に反する。残る七つが定理の値である。
逆に、七つの値がそれぞれ $m=1,6,11,45,51,209,660099$ で解になることは、代入して確かめられる。
\end{proof}

\begin{remark}
$13775$ が現れるのは、$d=2$ の初期元の $\rho^{-4}$ 倍が $f$ の零点になり、その第 $0$ 座標が $13778$ だからである。
本稿の証明は、この根の存在を予測する理由を与えるものではない。
\end{remark}

%======================================================================
\appendix
\section{Lean による形式化}\label{app:lean}
%======================================================================

最終定理は次の形をしている。$\N$ の割り算は切り捨てだが、$6\mid n(n+1)(n+2)$ なので、ここでは正確な商である。
\begin{verbatim}
theorem Tetrahedral.tetrahedral_classification (n : ℕ) :
    (∃ m : ℕ, n*(n+1)*(n+2)/6+1=m^2) ↔
      n=0 ∨ n=5 ∨ n=8 ∨ n=22 ∨ n=24 ∨ n=63 ∨ n=13775
\end{verbatim}

本文の各節と Lean のファイルの対応は表~\ref{tab:lean} のとおりである。
表のパスは、ディレクトリ \texttt{lean/Tetrahedral/} からの相対パスである。拡張子 \texttt{.lean} は省いた。

\begin{table}[htbp]
\centering
\small
\caption{本文と Lean の対応}\label{tab:lean}
\begin{tabular}{lL{11cm}}
\toprule
節 & ファイル\\
\midrule
\ref{sec:branches} & \texttt{Basic}, \texttt{Squarefree}, \texttt{IntegerParity}, \texttt{IntegerSquareObstructions}, \texttt{Branches}\\
\ref{sec:parameters} & \texttt{SpecialPell}, \texttt{QuadraticParameters}, \texttt{SelectedParameters}, \texttt{SquareFactor}, \texttt{D3Parameters}, \texttt{NormEleven/Parameters}\\
\ref{sec:norm11} & \texttt{Quartic/D*Ring}, \texttt{Quartic/D1}, \texttt{D2}, \texttt{D3}, \texttt{D11}, \texttt{D66}, \texttt{NormEleven/Certificates}\\
\ref{sec:division} & \texttt{Quartic/D1Division}, \texttt{Quartic/D2Division}, \texttt{Quartic/D3Division}, \texttt{Quartic/D11}, \texttt{Quartic/D66}\\
\ref{sec:units} & \texttt{Quartic/UnitGap}, \texttt{Quartic/RelativeUnitGap}, \texttt{Quartic/D*Trace}\\
\ref{sec:signs} & \texttt{PositiveOrbit}, \texttt{OrbitCongruence}, \texttt{NormEleven/PositiveOrbitData}, \texttt{NormEleven/D2Congruences}\\
\ref{sec:dyadic} & \texttt{Binomial}, \texttt{QuadraticPolynomial}, \texttt{FactoredExponent}, \texttt{NormDerivative}, \texttt{RelativeDirection}, \texttt{NormEleven/RelativeDirectionData}, \texttt{NormEleven/ParityDirectionData}\\
\ref{sec:exponents} & \texttt{OrbitReflection}, \texttt{FiniteDifference}, \texttt{NormEleven/DyadicSimple}, \texttt{NormEleven/DyadicCertificates}, \texttt{NormEleven/D2Entrance}, \texttt{NormEleven/D2Dyadic}, \texttt{NormEleven/D11Dyadic}, \texttt{NormEleven/D66Square}, \texttt{NormEleven/D66Dyadic}\\
\ref{sec:conclusion} & \texttt{NormEleven/D1}, \texttt{D2}, \texttt{D3}, \texttt{D11}, \texttt{D66}, \texttt{Classification}\\
\bottomrule
\end{tabular}
\end{table}

本文と Lean では、記号と名前が一部異なる。主なものを表~\ref{tab:names} に挙げる。

\begin{table}[htbp]
\centering
\small
\caption{本文と Lean の記号}\label{tab:names}
\begin{tabular}{ll}
\toprule
本文 & Lean\\
\midrule
$\omega=\sqrt e$ & \texttt{q}\\
$\theta$ & \texttt{s}（$d=66$ では \texttt{t}）\\
$\Nm(x)$ の二係数 & \texttt{relH x}, \texttt{relK x}\\
$\tau$ & \texttt{tau}\\
$\sigma$ & \texttt{sig}\\
基点 & centre\\
方向 $U$ & \texttt{step}\\
$h$ & width\\
$\chi$ & parity character\\
\bottomrule
\end{tabular}
\end{table}

次の三つのファイルは、\texttt{scripts/} の Python スクリプトが生成した証明である。
\begin{itemize}
\item \texttt{NormEleven/PositiveOrbitData}
\item \texttt{NormEleven/RelativeDirectionData}
\item \texttt{NormEleven/ParityDirectionData}
\end{itemize}
生成されたファイルも通常の Lean の証明であり、カーネルが検査する。
スクリプトは、同じファイルを再生成できることを確かめるために置いている。

\begin{thebibliography}{9}
\bibitem{GPZ94} J.~Gebel, A.~Peth\H{o}, H.~G.~Zimmer, Computing integral points on elliptic curves, \emph{Acta Arith.} \textbf{68} (1994), 171--192.
\bibitem{ST94} R.~J.~Stroeker, N.~Tzanakis, Solving elliptic diophantine equations by estimating linear forms in elliptic logarithms, \emph{Acta Arith.} \textbf{67} (1994), 177--196.
\bibitem{OEIS} The On-Line Encyclopedia of Integer Sequences, A000292 and A216269, \url{https://oeis.org}.
\end{thebibliography}

\end{document}
